% 第三部分 C：不依赖结果的章节（由 main_zh.tex 引入）。

\section{评判未完成的解答}\label{sec:prefix}

树状 rollout~\cite{hou2025treerl} 对一个解答采样多个前缀，续写其中一部分，并让各续写共享同一前缀。若评判模型能尽早看出哪些前缀会导向正确答案，就可以决定续写哪些前缀：省下续写死路的代价，把样本集中在策略能够成功的地方。我们先离线测量一个 4B 评判模型能否做出这种判断，再用它来训练。

\paragraph{前缀研究。} 我们使用 GRPO 的验证题（从训练中留出的 500 道 MATH 训练题，见第二部分第 2 节）。基座策略按 GRPO 的方式每题采样 4 个解答（温度 1，训练提示）。每个解答若足够长，就在 128、256 和 512 个词元处截断。每个前缀续写 4 次，直到训练长度（总共 1024 个词元），由验证器检查。一个前缀的续写中最终正确的比例就是它的蒙特卡洛价值。所比较的信号为：
\begin{itemize}
\item \kev{} 对 \emph{on track}（“这个未完成的解答是否正朝着正确的最终答案前进……？”）回答“是”的概率；
\item \kev{} 对第二部分 \emph{correct} 问题回答“是”的概率，把前缀当作所给解答；
\item AR 评判对 \emph{on track} 回答“是”的概率；
\item 策略自身在该前缀上的平均词元对数概率，不需要任何额外代价。
\end{itemize}
对每个截断长度，我们报告每种信号对“续写最终正确”的 AUROC：一是合并所有题目，二是在题目内部——在同一道题的前缀之间做选择时，重要的正是后者。我们还报告每种信号与蒙特卡洛价值的 Spearman 相关系数。

\paragraph{剪枝实验。} 在 4 个解答都超过截断长度的题目上，由一种信号保留它排得最高的 2 个前缀。我们把所保留前缀的平均价值与全部 4 个前缀的平均价值比较，后者即随机选择在期望上保留的价值。我们还与一个预言机比较：它按一半续写排序，在另一半续写上计分。区间为按题目做 bootstrap 的 95\% 区间。我们还评估一条固定阈值规则：当评判模型的 $p(\text{yes})<0.3$ 时剪掉该前缀。对这条规则，我们报告它剪掉多少前缀、被剪掉的前缀中有多少是死前缀（价值为 0），以及它丢掉了多少正确的续写。

\section{树状 GRPO}\label{sec:tree}

每个提示的 8 个叶子来自 $R=4$ 个采样出的 128 词元根前缀。其中 $K=2$ 个被续写，每个从精确的词元前缀出发采样 4 次。\texttt{tree\_kev} 保留 \kev{} \emph{on track} 分数最高的两个根；\texttt{tree\_random} 随机保留两个。在分布上，\texttt{tree\_random} 等价于只采样两个根。已经结束的根本身就是完整解答，若被保留，它的 4 个叶子就是它的副本。叶子的训练方式与第二部分的扁平解答相同：完整的根词元和续写词元、验证器的奖励，以及在该提示 8 个叶子上的组基线。其余一切沿用第二部分的 \texttt{rlvr} 设置。每隔十步，未被保留的根也会被续写，仅用于记录：这测量了该步上的选择相对于保留全部根、以及相对于事后最好的两个根的收益。

两种效应方向相反。评判模型判断正确时，按评判选根会提高正确叶子的比例，从而让 GRPO 得到更多有学习信号的组。但这也使保留下来的根不再是策略的样本，梯度也就不再是策略梯度的无偏估计。比较 \texttt{tree\_kev} 与 \texttt{tree\_random}（树形相同、选择无偏）可以分离出选择的作用；比较 \texttt{tree\_random} 与第二部分的 \texttt{rlvr} 可以分离出树形的作用。
